% Written audit and bridge, 11 October 2026. Input fragment.
\section{The fixed dictionary behind the actual signed rows}
\label{sec:fixed-dictionary-bridge}

This section concerns the local family whose parent character and ray
modulus are fixed before the scale. It does not assert a uniform density
estimate for arbitrary parent characters with growing conductor. That
stronger assertion is unnecessary here: the moving factors in both
transformations are retained in the complete integer row.

Write $m_0=2(1-\omega)$. Let $Q$ be the \emph{literal} fixed ideal used
by the source's fixed-inducing-row predicate. Its hypotheses include
$Q\subseteq(72)$ and the stated containments in the original fixed ray
moduli. The relative slot characters have presentations of modulus
dividing $Q$. These containments use the ideal convention that
$Q\subseteq\mathfrak m$ means $\mathfrak m\mid Q$.
For the actual initial family $F$, this ideal is
\[
 Q=\operatorname{internalQ}
       (F.\mathrm{modulus}\cap(72))\,\eta_0
   =F.\mathrm{modulus}\cap(72)\cap\mathfrak q_{\eta_0},
\]
as defined by \path{NaturalFixedRaySourceFixedIdeal.sourceFixedIdeal}
and \path{sourceExceptional}. Thus the needed containments belong to
the literal source family; they are not a later enlargement of its
exceptional set. In a transformed child this ideal remains fixed even
though the separate whole-product character can move.

\paragraph{A finite collection chosen before arithmetic labels.}
Take the primary ray characters modulo $12$, the six unit-row characters
modulo $36$, and the allowed relative slot characters, together with their
inverses. On ideals prime to $Q$ their products belong to the finite group
of ray characters modulo $Q$. Multiply these products by each of the
original fixed parent characters; the resulting finite collection is a
union of fixed translates of that group. The parent itself need not have
modulus dividing $Q$. Different zero extensions at primes dividing $Q$
introduce only finitely many additional presentations. Thus there is a finite list
$\{\eta_\lambda:\lambda\in\Lambda\}$ containing every product that will
occur, chosen before the common supports, frequencies, or rows. Products
over an arbitrarily large number of active primes do not enlarge this
list: the reciprocity phases are characters of a single fixed residue
group. They do not introduce one new independent character for each prime.
Keep distinct translates of the $Q$-character group separate when defining
zero bins. The finite union is used to make constants simultaneous; it
does not imply that eligibility for one translate makes all other
translates nonprincipal. Partitioning the fixed outer labels by their
parent translate costs a fixed finite factor in the eventual joint count.

The source realizes the primary phases by
\path{FixedRay.primaryRayCharacter}, whose modulus is $(12)$, and unit
phases by \path{HeckeUnitRows.character}, whose modulus is $(36)$.
For the relative factors, the exact source formula is
\[
 \nu_\theta=\chi_\theta\eta_0^{-1},\qquad
 Q\subseteq M,\ Q\subseteq\mathfrak q_{\eta_0}
 \ \Longrightarrow\ Q\subseteq\mathfrak q_{\nu_\theta}.
\]
This is \path{NaturalFixedRaySourceCharacters.relativeCharacter_fixedQ}.
The allocation theorem \path{AllocatedRayDictionary.originalImage_fixedQ}
preserves it for every surviving slot. The full source identities
\path{both_common_rectangles} retain the original relative characters,
profiles and independent heights; allocation changes which original slots
remain, not their character dictionary.

\paragraph{Exact coefficient identity.}
Use the signs and primary-generator conventions of the signed-row
reconstruction. For either first-side orientation the complete second
integer row is
\[
 w=u\,eF_{\pm}GVh\ne0,
\]
where the fixed unit $u$ may instead be absorbed into the finite list.
Every actual residual character coefficient, including a surviving
relative slot twist, is of the form
\[
 c_{\eta_\lambda}(I)\,\mathfrak s_{m_0^6w}(I)
       \mathbf1_{(I,\mathfrak R)=1}.                 \tag{FD}
\]
Here $\mathfrak R$ is the retained natural and common puncture. This is an
identity on all ideals, including the zero values. On supported ideals
it follows by evaluating the signed identities at the primary generator;
on unsupported ideals both sides vanish because of $m_0$. A canceled
exponent at an active prime produces a factor
$\mathbf1_{(I,\mathfrak p)=1}$ and therefore belongs to $\mathfrak R$.
The second extraction contributes $GVh$ multiplicatively to the same
integer row. The complete second coprimality projector is expanded with
its original signs. Extracting one of its assigned divisors $d$ uses
$\psi(dn)=\psi(d)\psi(n)$; it does not impose an additional condition
$(d,n)=1$. A zero outside coefficient annihilates that term. These facts
are exactly the distinctions made in the signed-row reconstruction.

The two first-side orientations can be treated separately, at a fixed
factor of two. Within the square of one first-side column the reflected
second rectangle has the row $-w$, which is covered by the same finite
unit-character list. There is no claim that $F_+$ and $F_-$ are equal.
If a whole rectangle is conjugated, retain that conjugation of the whole
character: its zero maximum is unchanged after reflection of the height
coordinate. Do not replace $\overline{\eta\,\mathfrak s_w}$ by
$\overline\eta\,\mathfrak s_w$. Such a replacement could change the
parent translate and its exceptional rows.

Choose the natural character $\chi_{\lambda,w}$ for the unmasked part of
(FD). The natural constructor gives
\[
 N\mathfrak q_{\chi_{\lambda,w}}
 \le N\mathfrak q_{\eta_\lambda}\,N(m_0)N(72)\,Nw
 \le K Nw,                                             \tag{FC}
\]
with one fixed $K$. Its ideal coefficients equal those of
\texttt{momentData $\eta_\lambda$}; the new theorem
\path{natural_row_moment_dictionary} verifies this equality.
One must use this natural representative for (FC), since the abstract
choice inside \path{RowData.character} has no stated conductor bound.
The primitive inducing character has still smaller conductor.

\paragraph{Literal exclusion survives the change of presentation.}
For a nonzero mask and row, the source proves
\[
 \operatorname{FI}_Q(\eta;m_0r,1,w)
 \ \Longleftrightarrow\
 \operatorname{FI}_Q(\eta;m_0,1,w),                       \tag{FE1}
\]
by \path{FixedRowMask.fixedInducingRow_mul_mask_iff}.
If $Q\subseteq\mathfrak q_\nu$, it also proves
\[
 \operatorname{FI}_Q(\eta\nu;m_0,1,w)
 \ \Longleftrightarrow\
 \operatorname{FI}_Q(\eta;m_0,1,w),                       \tag{FE2}
\]
by \path{FixedRayInducingTransport.fixedInducingRow_product_iff}.
The reverse implication is not obtained by canceling a possibly zero
character value. Multiplication by $\nu^{-1}$ recovers the original
character with the additional modulus mask of $\nu$; (FE1) then removes
that mask from the \emph{inducing predicate}. This is why the equivalence
remains valid on the literal zero-extended family.

Apply (FE1), (FE2), and (FD) to the actual source exclusion. They identify
it with the same predicate on $\chi_{\lambda,w}$. In particular every
allowed relative twist in this same parent translate is nonprincipal on
an eligible row. No conclusion is made here about another parent
translate. The source
states this consequence directly as
\path{NaturalFixedRaySourceCharacters.sector_nonprincipal}.
For the second common extraction,
\path{SecondExceptionalFamily.Family.inducing_iff} and
\path{paired_inducing_iff} implement precisely the reduced-row and
reflected-row versions of this transport. The ideal $Q$ is not enlarged.
It would be incorrect simply to replace the source exclusion by
nonprincipality of one preferred primitive representative.

\paragraph{Zero bins and the density input.}
For one fixed parent translate, take the maximum of the zero real parts
over its allowed $Q$-character twists, and use the buffered height boxes
of the source. Whole-character inverses have the same maximum and need
not be modeled by an additional translate. Each actual character in
that maximum is a masked presentation
of one of (FD). Finite Euler factors are nonzero in $\Re s>0$, so its
zeros in these boxes equal those of its primitive inducing character.
Removing sixth powers, $w=ub^6$, introduces further Euler holes and
does not change these positive-real zeros. The fixed-dictionary density
proposition therefore applies to the selected free core $u$, followed
by the height-uniform sixth-power lift.

If a bin is selected by the maximum over all allowed labels in this translate,
some label attains that maximum. The union over this \emph{fixed} finite
list costs only a constant. The upper side of the buffered bin bounds
every surviving prime-slot character, while the lower side supplies the
zero used for density. This step neither sums over the moving arithmetic
labels nor freezes them before counting. Those labels are recovered from
$w$ and the missing radical by the signed-row reconstruction.

Consequently, for $7/10\le\delta\le3/4$, the actual eligible integer rows
in the selected high bin satisfy
\[
 \#\{w:Nw\le H\}\ll_{\varepsilon}
       (1+T)^J H^{\rho(\delta)+\varepsilon},\qquad
 \rho(\delta)=1-\frac{10\delta}{3(5-2\delta)},
\]
up to the stated arbitrarily small bin-buffer loss. The constants and
threshold precede $T$, every live common-support label, and $H$.
Polynomial-size masks do not enter the dictionary size or the primitive
conductor in (FC); their smooth estimates are paid separately by uniform
Euler deletion. The pointwise upper bounds use the common buffered bin,
whereas the count only needs the label attaining its lower edge.

\paragraph{Fourier kernels whose parameter varies with labels.}
The common Fourier argument also permits a scale parameter $A_\ell$
depending on a retained label $\ell$. Uniform Euler derivative bounds
give Fourier densities satisfying
$|f_A(t)|\le C(1+|t|)^{-2r}$, uniformly in $A$. Choose once and for all
$d\mu(t)=C(1+|t|)^{-2r}\,dt$ and write
$f_{A_\ell}(t)\,dt=a_\ell(t)d\mu(t)$, where $|a_\ell(t)|\le1$.
The integral of the joint vector then contains the diagonal contraction
$\operatorname{diag}(a_\ell(t))$. Its operator norm is at most one,
so Minkowski's inequality gives exactly the same joint bound, with no
extra label count. This avoids the unjustified assertion that the
individual Fourier measures for distinct $A_\ell$ are equal.

This establishes the fixed-dictionary, conductor, mask, and literal
eligibility bridge from the actual signed-row identities. The remaining
assembly of the physical moment must still verify its normalization,
comparison and exceptional terms, and its use of the reserved exponent
margins. Those conclusions do not follow from this bridge alone.
